\documentclass[aspectratio=169]{beamer}

\usepackage[T1]{fontenc}
\usepackage[utf8]{inputenc}
\usepackage{lmodern}
\usepackage{amsmath}
\usepackage{booktabs}
\usepackage{array}
\usepackage{tabularx}
\usepackage{tikz}
\usepackage{pgfplots}
\usepackage{pgfplotstable}

\pgfplotsset{compat=1.18}

\definecolor{CorpBlue}{HTML}{0F4C81}
\definecolor{CorpTeal}{HTML}{1C7C7D}
\definecolor{CorpOrange}{HTML}{D17A22}
\definecolor{CorpLight}{HTML}{F3F6FA}
\definecolor{CorpGray}{HTML}{5B6573}

\usetheme{Madrid}
\setbeamercolor{structure}{fg=CorpBlue}
\setbeamercolor{palette primary}{bg=CorpBlue,fg=white}
\setbeamercolor{palette secondary}{bg=CorpTeal,fg=white}
\setbeamercolor{palette tertiary}{bg=CorpGray,fg=white}
\setbeamercolor{frametitle}{bg=CorpLight,fg=CorpBlue}
\setbeamercolor{block title}{bg=CorpBlue,fg=white}
\setbeamercolor{block body}{bg=CorpLight,fg=black}
\setbeamertemplate{navigation symbols}{}

\newcommand{\code}[1]{\texttt{#1}}

\title{NILM Signal Disaggregation\\Approcci sviluppati, benchmark}
\author{}
\date{}

\begin{document}

\begin{frame}
  \titlepage
\end{frame}

\begin{frame}{Executive summary}
\small
\begin{itemize}
  \item Il progetto affronta il task NILM: stimare i contributi dei singoli elettrodomestici a partire dal solo segnale aggregato \code{w\_medio} e da un inventario survey per IMEI.
  \item Sono stati sviluppati e testati approcci event-based, HMM, FHMM, template matching e una famiglia Gurobi con formulazioni binarie e multi-livello.
  \item Sul benchmark con copertura completa di 6 abitazioni idonee, \code{fhmm\_1\_dc} e' il miglior modello medio con MAE di ricostruzione pari a 27.64 W.
  \item \code{gurobi\_daywise} e' competitivo come piattaforma di sviluppo: MAE medio 77.02 W, piu' interpretabile nei vincoli e piu' estendibile verso regole operative esplicite.
  \item Le estensioni survey-aware e multi-livello sono gia' implementate; il passo successivo consigliato e' una campagna benchmark sistematica su tutte le 6 abitazioni.
\end{itemize}
\end{frame}

\begin{frame}{Obiettivo di business e perimetro}
\small
\begin{itemize}
  \item Obiettivo: trasformare una misura energetica aggregata per abitazione in una stima per dispositivo, senza sensori invasivi e senza labels temporali di ground truth.
  \item Vincoli reali: segnale rumoroso, dispositivi simultanei, carichi sempre accesi, profili d'uso eterogenei e questionari incompleti.
  \item Dataset operativo: 6 abitazioni idonee con storico sufficiente e inventario dispositivi associato a ogni IMEI.
  \item Output atteso per il business: CSV di disaggregazione, grafici temporali, report energetico e benchmark comparativo tra approcci.
\end{itemize}

\vspace{0.5em}
\begin{block}{Abitazioni idonee}
\small
H1 = 86684007269866, H2 = 86684007269887, H3 = 86684007269889, H4 = 86853106211162, H5 = 86853106211173, H6 = 86853106211179.
\end{block}
\end{frame}

\begin{frame}{Dati e pipeline end-to-end}
\small
\begin{columns}[T]
\column{0.52\textwidth}
\begin{block}{Input}
\begin{itemize}
  \item JSON IoT per IMEI con segnale \code{w\_medio} a granularita' minuto.
  \item Inventario survey legacy \code{device\_usage\_by\_imei}.
  \item Inventario survey V2 \code{device\_usage\_by\_imei\_v2} con bounds, finestre orarie e stagionalita'.
\end{itemize}
\end{block}

\column{0.48\textwidth}
\begin{block}{Pipeline}
\begin{enumerate}
  \item Preprocessing e ricampionamento a 1 minuto.
  \item Costruzione dei \code{DeviceProfile}.
  \item Esecuzione dell'approccio scelto tramite interfaccia comune \code{run(signal, devices)}.
  \item Salvataggio di \code{disaggregation.csv}, plot temporali e report energetico.
  \item Benchmark proxy basato su errore di ricostruzione, energia e consistenza temporale.
\end{enumerate}
\end{block}
\end{columns}

\vspace{0.3em}
\[
\widehat{w}(t) \approx \sum_i p_i(t), \qquad
\mathrm{MAE} = \frac{1}{T}\sum_t |w_{\mathrm{total}}(t) - \widehat{w}(t)|
\]
\end{frame}

\begin{frame}{Deep dive: \code{devices.py}}
\small
\begin{itemize}
  \item \code{DeviceProfile} e' il contratto condiviso da tutti i modelli. Contiene potenze minime, tipiche e massime, durata minima e tipica, duty cycle, frequenza, prior, flag always-on e metadati opzionali per vincoli survey-aware.
  \item \code{DEVICE\_KNOWLEDGE\_BASE} incapsula la knowledge base ingegneristica di 15 categorie: frigoriferi, carichi di cottura, microonde, climatizzazione, TV, computer, auto elettrica e altri.
  \item \code{get\_device\_profiles()} converte l'inventario legacy in profili operativi. La presenza del device imposta \code{prior\_weight = 1.0}; l'assenza lo degrada a 0.05 invece di escluderlo del tutto.
  \item Questo file e' il punto centrale di allineamento tra survey, priors, benchmark e approcci: cambia una sola volta e si aggiorna l'intera pipeline NILM.
\end{itemize}
\end{frame}

\begin{frame}{Deep dive: \code{device\_usage\_by\_imei\_v2}}
\small
\begin{itemize}
  \item La directory V2 deriva da \code{build\_questionnaire\_device\_usage\_v2.py}, che normalizza un export questionario grezzo in JSON strutturati per IMEI.
  \item Ogni device salva bounds numerici, non solo un valore medio: \code{usage\_frequency\_per\_week\_min/max}, \code{duration\_minutes\_min/max}, finestre di start, ore giornaliere, mesi attivi e warning di qualita'.
  \item \code{get\_device\_profiles\_v2()} riduce questi bounds a valori rappresentativi per l'inferenza, ma mantiene anche i limiti massimi utili nei modelli survey-aware.
  \item Questa struttura abilita euristiche piu' forti: commitment window guidata dal survey, cap di durata, numero massimo di blocchi, finestre orarie e stagionalita'.
\end{itemize}

\vspace{0.4em}
\begin{block}{Esempio pratico}
\small
Per \code{Lavatrice} si possono avere frequenza settimanale 1--2, durata 0--60 minuti e finestra 06:00--10:00; per \code{Climatizzatore} si possono avere 2--4 ore/giorno e mesi attivi 6--9.
\end{block}
\end{frame}

\begin{frame}{Approccio: Event-based}
\small
\begin{itemize}
  \item Modello piu' semplice: rileva eventi ON/OFF tramite \code{delta\_w = signal.diff()} e soglia fissa, poi associa ogni salto positivo al device con potenza tipica piu' vicina.
  \item Ogni evento ON viene accoppiato con il primo OFF compatibile entro \code{dur\_typical x 3}; in mancanza di OFF usa una durata di default.
  \item Punti di forza: spiegabile, veloce, utile come baseline interpretabile.
  \item Limiti: fragile con carichi sovrapposti o salti di potenza non puliti.
  \item Benchmark 6 case: MAE medio 151.44 W.
\end{itemize}
\end{frame}

\begin{frame}{Approccio: Event-based con prior}
\small
\begin{itemize}
  \item Mantiene la logica event-based ma sostituisce la sola prossimita' in potenza con uno score bayesiano: likelihood sul salto osservato per prior di presenza e frequenza d'uso dichiarata.
  \item L'inventario survey non elimina i device assenti, ma li penalizza in modo forte tramite \code{prior\_weight}.
  \item Punti di forza: migliore allineamento con il contesto reale rispetto al puro matching per potenza.
  \item Limiti: rimane dipendente dalla qualita' del rilevamento eventi e non modella bene i carichi simultanei lunghi.
  \item Benchmark 6 case: MAE medio 150.40 W, miglioramento marginale rispetto all'event-based puro.
\end{itemize}
\end{frame}

\begin{frame}{Approccio: HMM per singolo device}
\small
\begin{itemize}
  \item Per ogni device presente viene fittato un \code{GaussianHMM} a 2 stati ON/OFF sul residuo corrente, con inizializzazione delle medie da knowledge base.
  \item La sequenza di device viene stimata in cascata: dopo ogni device, il contributo stimato viene sottratto dal residuo globale.
  \item Punti di forza: struttura probabilistica chiara e uso di una libreria standard.
  \item Limiti: decomposizione sequenziale, forte sensibilita' agli errori iniziali e risultati deboli in presenza di simultaneita'.
  \item Benchmark 6 case: MAE medio 2110.03 W, il peggiore tra gli approcci a copertura completa.
\end{itemize}
\end{frame}

\begin{frame}{Approccio: FHMM semplificato}
\small
\begin{itemize}
  \item Modella tutti i device simultaneamente con catene latenti binarie indipendenti e inferenza greedy tipo coordinate-ascent.
  \item Ogni stato \code{x\_i(t)} viene aggiornato minimizzando l'errore di ricostruzione sul residuo generato dagli altri device; a convergenza si applica smoothing temporale sui blocchi troppo brevi.
  \item Punti di forza: rappresenta combinazioni simultanee senza esplosione combinatoria completa.
  \item Limiti: potenza binaria fissa e nessuna modellazione esplicita di carichi always-on.
  \item Benchmark 6 case: MAE medio 71.35 W, seconda miglior media tra i modelli con copertura completa.
\end{itemize}
\end{frame}

\begin{frame}{Approccio: \code{fhmm\_1} peak baseline}
\small
\begin{itemize}
  \item Estende FHMM introducendo una baseline costante per i carichi always-on e lavorando solo sul residuo positivo dei device event-driven.
  \item Usa una commitment window: una volta acceso, il device rimane ON per almeno \code{dur\_typical\_min}; i device bursty come il microonde mantengono potenza fissa.
  \item Punti di forza: piu' realistico del FHMM base per frigoriferi e cicli continui.
  \item Limiti: la baseline a potenza di picco puo' sovrastimare gli always-on in diversi contesti abitativi.
  \item Benchmark 6 case: MAE medio 80.92 W.
\end{itemize}
\end{frame}

\begin{frame}{Approccio: \code{fhmm\_1\_dc} duty-average baseline}
\small
\begin{itemize}
  \item Mantiene la struttura di \code{fhmm\_1} ma usa per gli always-on una baseline media pari a \code{p\_typical * duty\_cycle}, anziche' la potenza di picco.
  \item Questo riduce la sovra-assegnazione sistematica dei carichi frigoriferi e libera residuo utile per i carichi intermittenti.
  \item Punti di forza: ottimo compromesso tra semplicita', stabilita' e accuratezza.
  \item Limiti: resta un modello euristico, con potenza assegnata principalmente guidata dal residuo.
  \item Benchmark 6 case: MAE medio 27.64 W, miglior risultato globale tra gli approcci con copertura completa.
\end{itemize}
\end{frame}

\begin{frame}{Approccio: \code{fhmm\_1\_survey}}
\small
\begin{itemize}
  \item Variante survey-aware di \code{fhmm\_1}: la soglia di attivazione viene modulata in base a finestre di start, mesi attivi, durata massima, frequenza massima e ore giornaliere del questionario V2.
  \item Introduce euristiche operative molto vicine a un ragionamento di dominio: ad esempio limita i blocchi giornalieri o favorisce l'accensione dentro la finestra oraria dichiarata.
  \item Punti di forza: piu' adatto a dispositivi con pattern temporali forti e questionario ricco.
  \item Limiti: benchmark ancora parziale, dipendente dalla qualita' del survey V2 disponibile per la casa.
  \item Benchmark attuale: MAE medio 49.50 W su 2 abitazioni.
\end{itemize}
\end{frame}

\begin{frame}{Approccio: Template matching}
\small
\begin{itemize}
  \item Costruisce per ogni device un template rettangolare di potenza tipica e durata tipica, poi calcola cross-correlazione normalizzata con il segnale aggregato.
  \item I picchi sopra soglia vengono assegnati in modo greedy senza overlap, con priorita' ai device piu' potenti.
  \item Punti di forza: semplice da interpretare e rapido da eseguire.
  \item Limiti: i consumi reali raramente seguono template rettangolari puliti; la metodologia degrada molto con rumore e device sovrapposti.
  \item Benchmark 6 case: MAE medio 1149.56 W.
\end{itemize}
\end{frame}

\begin{frame}{Approccio: Gurobi weekly MIQP}
\small
\begin{itemize}
  \item Formula il NILM come problema di Mixed-Integer Quadratic Programming: minimizza l'errore quadratico tra segnale aggregato e somma dei contributi binari dei device su chunk settimanali.
  \item L'architettura supporta variabili di stato, attivazioni, spegnimenti, vincoli di durata e penalita' sulle accensioni, con forte controllabilita' modellistica.
  \item Punti di forza: massima espressivita' tra gli approcci attuali e ottimo fit con vincoli di business espliciti.
  \item Limiti: costo computazionale elevato e benchmark ancora molto parziale.
  \item Benchmark attuale: MAE medio 147.01 W su 1 sola abitazione.
\end{itemize}
\end{frame}

\begin{frame}{Approccio: Gurobi daywise}
\small
\begin{itemize}
  \item Variante piu' pragmatica: assegna una baseline sempre accesa ai frigoriferi, la sottrae dal segnale e risolve il residuo giorno per giorno con un solver binario Gurobi piu' piccolo.
  \item In \code{gurobi\_methods.py} sono state aggiunte anche le estensioni \code{constrained\_v1} e \code{constrained\_v2}, che introducono variabili di transizione, durata massima consecutiva e finestre orarie.
  \item Punti di forza: migliore trattabilita', architettura modulare e forte potenziale di controllo dei falsi positivi.
  \item Limiti: sensibilita' elevata alle abitazioni piu' difficili; H6 rimane un caso critico con MAE 212.60 W.
  \item Benchmark 6 case: MAE medio 77.02 W.
\end{itemize}
\end{frame}

\begin{frame}{Approccio: Gurobi multilevel 3 e 5 stati}
\small
\begin{itemize}
  \item Il nuovo modulo \code{gurobi\_multiple.py} generalizza il modello binario a piu' livelli di potenza: ogni device puo' essere OFF o scegliere uno tra 3 o 5 livelli discreti.
  \item I livelli sono ottenuti per interpolazione lineare tra \code{p\_min}, \code{p\_typical} e \code{p\_max}; il livello centrale coincide con la potenza tipica, gli altri si avvicinano ai bound senza coincidere con essi.
  \item Le varianti disponibili sono: unconstrained, \code{v1} con variabili \code{u/dw} e durata massima, e \code{v2} con finestre orarie per microonde e lavatrice.
  \item Punti di forza: rappresenta meglio device con potenza non binaria e apre la strada a modelli piu' realistici per forno, climatizzazione e lavatrice.
  \item Stato: implementato e integrato nel runner, ma benchmark sistematico ancora da completare.
\end{itemize}
\end{frame}

\begin{frame}{Focus Gurobi: evoluzione della famiglia di modelli}
\scriptsize
\begin{tabularx}{\textwidth}{>{\bfseries}p{0.20\textwidth}p{0.34\textwidth}p{0.12\textwidth}X}
\toprule
Versione & Idea chiave & Copertura & Stato \\
\midrule
Weekly MIQP & Solve settimanale binario con alta espressivita' modellistica & 1 casa & Benchmark iniziale \\
Daywise binary & Baseline always-on + solve giornaliero sul residuo & 6 case & Variante piu' matura \\
\code{v1} & Variabili \code{u/dw} e durata massima consecutiva & 0 case complete & Implementata, test da fare \\
\code{v2} & V1 + finestre orarie per device selezionati & 0 case complete & Implementata, test da fare \\
Multilevel 3/5 & Livelli di potenza multipli con versioni unconstrained/v1/v2 & 0 case complete & Pronta per benchmark \\
\bottomrule
\end{tabularx}

\vspace{0.4em}
\begin{block}{Messaggio chiave}
\scriptsize
La famiglia Gurobi non e' ancora la migliore in media rispetto a \code{fhmm\_1\_dc}, ma e' la piu' promettente come piattaforma industriale per introdurre vincoli espliciti, multi-stato e regole di dominio controllabili.
\end{block}
\end{frame}

\begin{frame}{Metodologia benchmark}
\small
\begin{itemize}
  \item Non esistono labels temporali di ground truth; per questo il benchmark e' proxy e misura la qualita' di ricostruzione del segnale aggregato.
  \item La metrica principale usata in questa presentazione e' il \textbf{MAE di ricostruzione}, calcolato come media del valore assoluto del \code{residuo} salvato in ogni \code{disaggregation.csv}.
  \item Sono presenti anche altre metriche nel modulo \code{benchmark.py}: RMSE, errore energetico percentuale, residuo medio, numero di device trovati e consistenza temporale.
  \item Per evitare confronti distorti, il ranking principale considera solo gli approcci con copertura completa su tutte le 6 abitazioni idonee.
\end{itemize}

\vspace{0.4em}
\[
\mathrm{MAE}_{\mathrm{recon}} = \frac{1}{T}\sum_{t=1}^{T}
\left|w_{\mathrm{total}}(t) - \sum_i \widehat{w}_i(t)\right|
\]
\end{frame}

\begin{frame}{Confronto risultati: MAE medio su 6 case}
\small
\begin{center}
\begin{tikzpicture}
\begin{axis}[
    width=0.98\textwidth,
    height=0.63\textheight,
    ybar,
    bar width=9pt,
    ymin=0,
    ymax=2300,
    ylabel={Mean MAE reconstruction [W]},
    xtick=data,
  xticklabels={FHMM-1 duty avg,FHMM,Gurobi daywise,FHMM-1 peak,Event + prior,Event-based,Template,HMM},
    x tick label style={rotate=35,anchor=east,font=\scriptsize},
    ymajorgrids=true,
    grid style={dashed,gray!30},
    nodes near coords,
    every node near coord/.append style={font=\tiny, rotate=90, anchor=west},
]
\addplot[fill=CorpTeal, draw=CorpBlue] table[x expr=\coordindex, y=mean_mae_recon, col sep=comma] {data/mae_summary_full_coverage_ranked.csv};
\end{axis}
\end{tikzpicture}
\end{center}

\vspace{-0.2em}
\tiny
Lettura del grafico: \code{fhmm\_1\_dc} e' il best performer medio; \code{gurobi\_daywise} e' nel gruppo competitivo ma non ancora al vertice; HMM e template sono nettamente fuori scala.
\end{frame}

\begin{frame}{Confronto risultati: MAE per abitazione}
\scriptsize
\pgfplotstabletypeset[
    col sep=comma,
    string type,
    columns={label,H1,H2,H3,H4,H5,H6},
    columns/label/.style={column name=Approach, column type=l},
    columns/H1/.style={column name=H1, column type=c},
    columns/H2/.style={column name=H2, column type=c},
    columns/H3/.style={column name=H3, column type=c},
    columns/H4/.style={column name=H4, column type=c},
    columns/H5/.style={column name=H5, column type=c},
    columns/H6/.style={column name=H6, column type=c},
    every head row/.style={before row=\toprule, after row=\midrule},
    every last row/.style={after row=\bottomrule},
]{data/mae_matrix_full_coverage_ranked_short.csv}

\vspace{0.5em}
\tiny
Mappatura: H1=86684007269866, H2=86684007269887, H3=86684007269889, H4=86853106211162, H5=86853106211173, H6=86853106211179.
\end{frame}

\begin{frame}{Risultati parziali e stato sperimentale}
\small
\begin{columns}[T]
\column{0.50\textwidth}
\begin{block}{Approcci con benchmark parziale}
\begin{itemize}
  \item \code{fhmm\_1\_survey}: MAE medio 49.50 W su 2 case. Segnale incoraggiante, ma copertura ancora troppo limitata per ranking finale.
  \item \code{gurobi}: MAE 147.01 W su 1 casa. Utile come prova di formulazione, non ancora sufficiente come baseline industriale.
\end{itemize}
\end{block}

\column{0.50\textwidth}
\begin{block}{Approcci implementati, benchmark pending}
\begin{itemize}
  \item \code{constrained\_v1}: transizioni \code{u/dw} + durata massima consecutiva.
  \item \code{constrained\_v2}: v1 + finestre orarie per microonde e lavatrice.
  \item \code{gurobi\_multiple}: livelli di potenza 3 e 5 con versioni unconstrained, v1 e v2.
\end{itemize}
\end{block}
\end{columns}

\vspace{0.4em}
\begin{block}{Interpretazione manageriale}
\small
Le estensioni piu' recenti non sono ancora da considerare perdenti o vincenti: il loro valore attuale e' architetturale, perche' aprono un percorso chiaro verso modelli controllabili e piu' fedeli al comportamento reale dei device.
\end{block}
\end{frame}

\begin{frame}{Raccomandazioni operative}
\small
\begin{itemize}
  \item Produzione a breve termine: usare \code{fhmm\_1\_dc} come baseline migliore per accuratezza media e robustezza sulle 6 case benchmarkate.
  \item Ricerca applicata: mantenere \code{gurobi\_daywise} come ramo di sviluppo principale quando il requisito e' inserire vincoli espliciti e controlli per device specifici.
  \item Survey-aware roadmap: ampliare la copertura di \code{device\_usage\_by\_imei\_v2} per portare \code{fhmm\_1\_survey} e \code{gurobi v2} su tutte le abitazioni.
  \item Benchmark next step: lanciare in batch le 6 varianti multi-livello \code{gurobi\_multiple\_*} su tutte le 6 case e confrontarle con \code{fhmm\_1\_dc} e \code{gurobi\_daywise}.
  \item KPI di decisione: oltre al MAE, monitorare energia residua, falsi positivi su microonde/lavatrice e stabilita' temporale dei blocchi ON.
\end{itemize}
\end{frame}

\begin{frame}{Conclusioni}
\small
\begin{itemize}
  \item La codebase oggi copre un ventaglio ampio di approcci, dalle euristiche leggere ai modelli con ottimizzazione matematica esplicita.
  \item Il cuore informativo del progetto e' la combinazione tra knowledge base ingegneristica in \code{devices.py} e arricchimento survey in \code{device\_usage\_by\_imei\_v2}.
  \item Il benchmark attuale indica una gerarchia chiara: \code{fhmm\_1\_dc} come best performer medio, \code{FHMM} e \code{gurobi\_daywise} come seconde linee competitive, HMM e template non raccomandati.
  \item La famiglia Gurobi rimane il percorso piu' promettente per introdurre regole operative, multi-stato e explainability di tipo industriale.
\end{itemize}

\vfill
\centering
\textbf{Messaggio finale:} la piattaforma e' gia' adatta a una fase aziendale di valutazione comparativa e prioritizzazione di una roadmap product-oriented.
\end{frame}

\end{document}